%
% 22-fortsetzung.tex -- Analytische Fortsetzung entlang einer Kurve
%
% (c) 2026 Prof Dr Andreas Müller
%
\subsection{Analytische Fortsetzung entlang einer Kurve
\label{buch:analytisch:fortsetzung:subsection:kurve}}
Sei $\gamma$ eine Kurve in der komplexen Ebene und $f$ eine
Funktion, die in einer Umgebung des Punktes $\gamma(0)$ der
Kurve holomorph ist.
Die Funktion ist also in einem Kreis vom Radius $\varrho(\gamma(0))$
um den Punkt $\gamma(0)$ analytisch.
Nach \eqref{buch:analytisch:fortsetzung:eqn:taylor01}
kann die Funktion $f(z)$ in jedem Punkt $\gamma(t)$ wieder in eine
Potenzreihe mit Konvergenzradius $\varrho(\gamma(t))$ entwickelt
werden, so dass die beiden Potenzreihen auf dem Schnittgebiet der
beiden Konvergenzkreise übereinstimmen.
Für einen Punkt $\gamma(t)$ nahe dem Rand des Konvergenzkreises 
der Potenzreihe im Punkt $\gamma(0)$ kann der Konvergenzkreis
der neuen Reihe über den Rand des alten Konvergenzkreises hinausragen.
Die neue Potenzreihe erweitert also die Funktion über den ursprünglichen
Konvergenzkreis hinaus.
Indem dies entlang der Kurve immer wieder wiederholt wird, kann die
Funktion auf ein Umgebung der Kurve ausgedehnt werden, wie dies zum
für die Wurzelfunktion in
Abbildung~\ref{buch:analytisch:fortsetzung:fig:wurzel2}
illustriert wird.

\begin{beispiel}
Die Funktion $f(z)=\frac{1}{z}$ ist holomorph in $\mathbb{C}\setminus\{0\}$.
Sie $\gamma(t)=e^{2\pi it}$ der Einheitskreis.
Im Punkt $\gamma(0)=1$ kann die Funktion durch die geometrische Reihe
\begin{equation}
\frac{1}{z}
=
\frac{1}{1+(z-1)}
=
\sum_{k=0}^\infty (z-1)^k
\label{buch:analytisch:fortsetzung:bsp:geom1}
\end{equation}
dargestellt werden, die in $B_1(1)$ konvergent ist.
Der Punkt $\gamma(\frac{\pi}4)$ liegt innerhalb dieses Konvergenzkreises.
Die Entwicklung
\begin{equation}
\frac{1}{z}
=
\frac{1}{a+(z-a)}
=
\frac1a \frac{1}{1+\displaystyle\frac{z-a\mathstrut}{a}}
=
\sum_{k=0}^\infty
\frac{1}{a^k}(z-a)^k
\label{buch:analytisch:fortsetzung:bsp:geoma}
\end{equation}
ist eine Potenzreihenentwicklung im Punkt $a$ mit Konvergenzradius $|a|$.
Für $a=\gamma(\frac{\pi}4)$ ist die Reihe von
\eqref{buch:analytisch:fortsetzung:bsp:geoma},
die in einer Umgebung von $a$ mit der Reihe
\eqref{buch:analytisch:fortsetzung:bsp:geom1}
übereinstimmt.
Der Punkt $i\in\mathbb{C}$ ist nicht im Konvergenzkreis der Reihe 
\eqref{buch:analytisch:fortsetzung:bsp:geom1},
der Wert $f(i)$ kann also damit nicht berechnet werden,
wohl aber mit
\eqref{buch:analytisch:fortsetzung:bsp:geoma}.
\end{beispiel}

Im Beispiel führt die analytische Fortsetzung entlang der Kurve 
nach einem Umlauf wieder zur selben Entwicklung
\eqref{buch:analytisch:fortsetzung:bsp:geom1}
im Punkt $1$.
Dies ist jedoch nicht selbstverständlich, wie die Beispiele
im kommenden Abschnitt zeigen werden.


\begin{satz}
Seien $f(z)$ und $g(z)$ holomorphe Funktionen auf Gebieten $U$
bzw.~$V$ die mitsamt aller ihrer Ableitungen in einem Punkt
von $U\cap V$ übereinstimmen.
Dann stimmen $f(z)$ und $g(z)$ auf jedem wegzusammenhängenden
Gebiet, welches die Kurve $\gamma$ enthält, überein.
\end{satz}

\begin{proof}
Sei $z_0$ der Punkte, in dem $f(z)$ und $g(z)$ samt ihren
Ableitungen übereinstimmen.
In diesem Punkt haben $f(z)$ und $g(z)$ die gleiche Potenzreihenentwicklung.
Wir müssen zeigen, dass die Funktionen $f(z)$ und $g(z)$ auch ein einem
Punkte $z_1$ übereinstimmen, der in einem einfach zusammenhängenden Gebiet
$W\subset U\cap V$ enthalten ist.
Wir konstruieren zu diesem Zweck einen differenzierbaren Weg $\gamma$
von $z_0$ nach $z_1$ in $W$.
Die Überlegungen von
Abschnitt~\ref{buch:analytisch:analytisch:subsection:potenzreihe}
zeigen dann, dass die beiden Funktionen samt ihren Ableitungen entlang
dieses Weges ebenfalls übereinstimmen müssen.
\end{proof}

Bei der Konstruktion einer analytischen Fortsetzung muss man sich also
nicht auf die Verwendung von Potenzreihen beschränken.
Dies wird auch vom folgenden Beispiel illustriert.

\begin{beispiel}
Die Funktionen
\begin{equation*}
\renewcommand{\arraycolsep}{2pt}
\begin{array}{rcrclcrcl}
f
&\colon&
\displaystyle
\biggl\{ re^{i\varphi}\mid r > 0\wedge -\frac{\pi}2<\varphi<\frac{3\pi}2 \biggr\}
&\to&
\mathbb{C}
&:&
z=re^{i\varphi}
&\mapsto&
\sqrt{r}e^{i\varphi/2}
\\[7pt]
g
&\colon&
\displaystyle
\biggl\{ re^{i\varphi}\mid r > 0\wedge -\frac{3\pi}2<\varphi<\frac{\pi}2 \biggr\}
&\to&
\mathbb{C}
&:&
z=re^{i\varphi}
&\mapsto&
\sqrt{r}e^{i\varphi/2}
\end{array}
\end{equation*}
sind so definiert sein, dass das Argument $\varphi$ auf den jeweiligen
Definitionsgebieten stetig ist.
Beide Funktionen haben die Eigenschaft
\[
f(z)^2
=
re^{i\varphi} = z
=
g(z)^2,
\]
sie sind also Wurzelfunktionen.

Die Definitionsgebiete bestehen jeweils aus der komplexen Ebene, aus der
ein Teil der imaginären Achse entfernt worden ist.
Konkret ist
\[
U = \mathbb{C}\setminus\{-iy\mid y\ge 0\}
\qquad\text{und}\qquad
V = \mathbb{C}\setminus\{+iy\mid y\ge 0\}.
\]
Die Schnittmenge ist 
\[
U\cap V
=
\{x+iy\in\mathbb{C}\mid x < 0\}
\cup
\{x+iy\in\mathbb{C}\mid x > 0\}.
\]
An der Stelle $z=1$ ist $r=1$ und $\varphi=0$ für beide Funktionen,
die Funktionswerte und die Ableitungen stimmen also überein.
Es folgt, dass die beiden Funktionen für $\operatorname{Re}z>0$
und damit auch auf jedem wegzusammenhängenden Teilgebiet davon
übereinstimmen werden.

Andererseits gilt für Punkte auf der negativen reellen Achse
\[
\renewcommand{\arraycolsep}{2pt}
\begin{array}{rclclcr}
f(-x) &=& f(xe^{i\pi})  &=& \!\sqrt{x}e^{i\pi/2} &=& i\!\sqrt{x},\\
g(-x) &=& g(xe^{-i\pi}) &=& \!\sqrt{x}e^{i\pi/2} &=& -i\!\sqrt{x}.
\end{array}
\]
Es gibt aber auch kein wegzusammenhängendes Teilgebiet von $U\cap V$,
welches den Punkt $z=1$ und einen Teil der negativen reellen Achse
enthält.
\end{beispiel}

Auch in diesem Beispiel ist des die Singularität der Wurzelfunktion
an der Stelle $0$, die eine eindeutige Definition verhindert.
Der Satz zeigt aber noch etwas mehr.
Er garantiert nämlich, dass man eine Funktion eindeutig auf ein
grösseres Gebiet ausdehnen kann, wenn man sicherstellen kann,
dass es nur einen Weg zu einem der neuen Punkte gibt.

